\section{Full reconstruction results}
\label{sec:s:results}

\paragraph{Three optical designs and four reconstruction methods.}
\Cref{tab:s:e2e} reports W1, centroid error, PSNR and foreground SSIM from the same 3,600 reconstructions: three plates, four methods, 60 crops and five draws. The feed-forward decoder is fine-tuned separately for each plate, whereas TV and the PnP networks are independent of optical training. \Cref{tab:s:e2eintervals} gives paired cell-bootstrap intervals for W1. Negative changes favor the first named design. The end-to-end plate has a measurable advantage with its own decoder and under both PnP methods, while our plate has lower W1 under TV. The different directions preclude a universal ranking of the two learned plates.

\begin{table*}[t]
\centering\small
\widetab{\begin{tabular}{llrrrr}
\toprule
Reconstruction & Plate & W1 ($\mu$m) & Centroid ($\mu$m) & PSNR & FG-SSIM \\
\midrule
TV & Untrained & 0.371 & 0.259 & 36.70 & 0.850 \\
TV & Ours & 0.218 & 0.150 & 38.29 & 0.910 \\
TV & End-to-end & 0.222 & 0.151 & 38.07 & 0.905 \\
\midrule
Pretrained PnP & Untrained & 0.579 & 0.490 & 35.72 & 0.806 \\
Pretrained PnP & Ours & 0.290 & 0.224 & 37.56 & 0.900 \\
Pretrained PnP & End-to-end & 0.281 & 0.212 & 37.52 & 0.897 \\
\midrule
Neuron-trained PnP & Untrained & 0.271 & 0.197 & 38.01 & 0.902 \\
Neuron-trained PnP & Ours & 0.195 & 0.139 & 39.60 & 0.934 \\
Neuron-trained PnP & End-to-end & 0.189 & 0.131 & 39.70 & 0.936 \\
\midrule
Feed-forward & Untrained & 0.366 & 0.282 & 36.67 & 0.872 \\
Feed-forward & Ours & 0.265 & 0.196 & 38.04 & 0.912 \\
Feed-forward & End-to-end & 0.234 & 0.167 & 38.53 & 0.922 \\
\bottomrule
\end{tabular}
}
\caption{Three-plate comparison at $C=8$ and matched nominal dose, with 60 crops and five draws. Errors are in $\mu$m. All metrics in a row use the same reconstructions. Feed-forward networks are fine-tuned for their respective plates; their earlier training histories are specified in \cref{sec:s:joint}.}
\label{tab:s:e2e}
\end{table*}

\begin{table*}[t]
\centering\small
\begin{tabular}{lrr}
\toprule
Reconstruction & Ours vs. untrained & End-to-end vs. ours \\
\midrule
TV & $-41.3\%$ [$-44.5$, $-37.7$] & $+1.9\%$ [$+1.1$, $+2.7$] \\
Pretrained PnP & $-49.9\%$ [$-52.8$, $-46.6$] & $-3.0\%$ [$-5.0$, $-1.0$] \\
Neuron-trained PnP & $-27.9\%$ [$-29.5$, $-26.0$] & $-3.2\%$ [$-5.6$, $-1.2$] \\
Feed-forward & $-27.5\%$ [$-31.2$, $-24.7$] & $-11.6\%$ [$-14.4$, $-9.0$] \\
\bottomrule
\end{tabular}

\caption{Paired relative W1 changes (95\,\% cell-bootstrap intervals) for the same 60-crop comparison. Each change is the ratio of crop-mean errors minus one, after averaging the five draws.}
\label{tab:s:e2eintervals}
\end{table*}

\paragraph{Optical sampling repeats.}
\Cref{tab:s:main} adds our second optical sampling seed to the primary TV/PnP comparison, using the same untrained initialization. Each learned plate lowers W1 on every crop under each method. Seed 0 improves centroid error on 59 of 60 crops for neuron-trained PnP and all crops for the other methods. Seed 1 improves centroid error on all crops. These repeats test sampling variability, not independence of phase initialization.

\begin{table*}[t]
\centering\small
\widetab{\begin{tabular}{lrrrrl}
\toprule
Reconstruction & W1 ($\mu$m) & Centroid ($\mu$m) & PSNR & FG-SSIM & W1 change [95\% CI] \\
\midrule
TV & 0.218 & 0.149 & 38.31 & 0.910 & $-41.3\%$ [$-44.6$, $-37.5$] \\
Pretrained PnP & 0.289 & 0.222 & 37.58 & 0.900 & $-50.1\%$ [$-53.1$, $-46.7$] \\
Neuron-trained PnP & 0.194 & 0.136 & 39.64 & 0.934 & $-28.5\%$ [$-30.1$, $-27.1$] \\
\bottomrule
\end{tabular}
}
\caption{Our second optical sampling seed at $C=8$, 60 crops and five draws per crop, with matched nominal dose. Changes compare it with the shared untrained initialization; brackets are cell-bootstrap 95\,\% intervals. Seed 0 is in \cref{tab:s:e2e}.}
\label{tab:s:main}
\end{table*}

\begin{table*}[t]
\centering\small
\widetab{\begin{tabular}{@{}lllccc@{}}
\toprule
Reconstruction & Seed & Crops & W1 change [95\,\% CI] & W1 improved & Centroid change [95\,\% CI]\\
\midrule
TV & 0 & 12 prespecified & $-38$\,\% [$-43$, $-34$] & 12/12 & $-37$\,\% [$-41$, $-34$]\\
TV & 0 & 48 random & $-42$\,\% [$-46$, $-37$] & 48/48 & $-43$\,\% [$-49$, $-38$]\\
TV & 1 & 12 prespecified & $-39$\,\% [$-42$, $-36$] & 12/12 & $-38$\,\% [$-40$, $-36$]\\
TV & 1 & 48 random & $-42$\,\% [$-46$, $-37$] & 48/48 & $-44$\,\% [$-49$, $-38$]\\
Pretrained PnP & 0 & 12 prespecified & $-44$\,\% [$-50$, $-39$] & 12/12 & $-46$\,\% [$-53$, $-41$]\\
Pretrained PnP & 0 & 48 random & $-51$\,\% [$-54$, $-47$] & 48/48 & $-56$\,\% [$-59$, $-52$]\\
Pretrained PnP & 1 & 12 prespecified & $-44$\,\% [$-49$, $-38$] & 12/12 & $-46$\,\% [$-52$, $-40$]\\
Pretrained PnP & 1 & 48 random & $-51$\,\% [$-55$, $-47$] & 48/48 & $-57$\,\% [$-60$, $-52$]\\
Neuron-trained PnP & 0 & 12 prespecified & $-27$\,\% [$-34$, $-20$] & 12/12 & $-26$\,\% [$-35$, $-18$]\\
Neuron-trained PnP & 0 & 48 random & $-28$\,\% [$-30$, $-26$] & 48/48 & $-31$\,\% [$-34$, $-28$]\\
Neuron-trained PnP & 1 & 12 prespecified & $-29$\,\% [$-34$, $-25$] & 12/12 & $-29$\,\% [$-35$, $-24$]\\
Neuron-trained PnP & 1 & 48 random & $-28$\,\% [$-30$, $-27$] & 48/48 & $-31$\,\% [$-34$, $-29$]\\
\bottomrule
\end{tabular}
}
\caption{Separate evaluation on the prespecified and randomly sampled crops. Negative changes favor the learned plate. Both sets contain crops from the same six cells; the 48 additional crops expand spatial coverage rather than biological replication.}
\label{tab:s:extended}
\end{table*}

\paragraph{Reference-based and random crops.}
\Cref{tab:s:extended} separates the 12 prespecified crops from the 48 later random crops. All methods retain positive optical gains on the latter, with parameters already fixed. This addresses dependence on the initial structure-score selection, while preserving the cell and animal limitations of the dataset.

\paragraph{Comparison with untrained encoders.}
The separate centroid study in \cref{tab:s:untrained} uses pretrained PnP, the 12 prespecified crops and the expanded validation grid. Its phase, amplitude and DMD rows share the same $2\times$ aperture, relay, photon model, direction draws and tuning procedure. Learned centroid errors of $0.230$--$0.234\um$ are below all three random phase plates ($0.422$--$0.438\um$) and amplitude masks ($0.396$--$0.523\um$). The learned seed-0 plate reduces centroid error by 42--56\,\% against the amplitude masks and 45--47\,\% against the random phase plates. This shows that its initialization is representative of these controls. It does not show that learning succeeds from other initializations.

\begin{table*}[t]
\centering\small
\begin{tabular}{@{}lccc@{}}
\toprule
Encoder & Centroid ($\mu$m) & PSNR (dB) & FG-SSIM\\
\midrule
Learned, seed 0 & 0.230 & 36.15 & 0.868\\
Learned, seed 1 & 0.234 & 36.21 & 0.869\\
Random phase 0 & 0.431 & 34.81 & 0.784\\
Random phase 1 & 0.422 & 34.93 & 0.787\\
Random phase 2 & 0.438 & 34.85 & 0.782\\
Amplitude mask 0 & 0.523 & 33.70 & 0.703\\
Amplitude mask 1 & 0.419 & 34.37 & 0.766\\
Amplitude mask 2 & 0.396 & 34.36 & 0.763\\
DMD set 0 & 2.169 & 28.83 & 0.482\\
DMD set 1 & 2.169 & 28.81 & 0.479\\
Plate-free, $4\times$ aperture & 1.102 & 30.18 & 0.494\\
\bottomrule
\end{tabular}

\caption{Untrained-encoder control study at $C=8$: 12 prespecified crops, pretrained PnP. W1 was not computed in this study. The plate-free row uses a $4\times$ aperture in a separate study; its own random-plate control gives $0.446\um$/34.82\,dB. Other rows use $2\times$ apertures. Do not interpret the plate-free row as an exactly matched cross-aperture contrast.}
\label{tab:s:untrained}
\end{table*}

The two in-focus DMD pattern sets remain close in centroid error and image quality. Their low depth discrimination concerns this fixed-focus control, not all amplitude encoders. Indeed, the defocused amplitude masks recover depth and can match random phase plates in centroid error. The plate-free scan also retains weak finite-aperture diversity, so its observations are not an ideal zero-depth-information system.

\paragraph{Point-emitter diagnostics.}
For 25 interior lateral locations and 16 depths, we compare measurement columns under a Gaussian approximation of shot noise with unknown emitter brightness. Random phase plates assign the wrong depth in 1.5--1.7\,\% of cases with 32 directions; amplitude masks give about 5\,\%, the in-focus DMD about 25\,\%, and the plate-free control about 89\,\%. These are model-based emitter diagnostics, not measured resolution of dense structures. The plate-free diagnostic uses the original $2\times$ aperture; its neuronal control in \cref{tab:s:untrained} uses $4\times$.

\paragraph{Comparisons across priors.}
A descriptive pairing of learned plate + TV against random plate + neuron-trained PnP gives 19.6\,\% lower W1 (95\,\% CI: 16.4--22.5\,\%), improving 56 of 60 crops. It does not establish a general dominance relation: learned plate + pretrained PnP has 7\,\% higher mean W1 than random plate + neuron-trained PnP (interval $-2.6$ to $+17.3$\,\%). Optical and prior choices interact; the comparisons in this paragraph concern only the untrained plate and our design, not the end-to-end plate.
